\documentclass[a4paper]{article}
\usepackage[swedish]{babel}
\usepackage{tikz}
\usepackage{fontspec}
\usepackage{xcolor}
\usepackage{enumitem}
\usepackage{background}
\usepackage{coffee}

\usetikzlibrary{calc}

%\setmainfont{Liberation Serif}
\newfontfamily\calfamily{EmbossedBlackWide}
\pagestyle{empty}
\renewcommand{\section}[1]{\begin{center}§ #1.\end{center}}


\SetBgScale{1}
\SetBgAngle{0}
\SetBgPosition{0,0}
\SetBgColor{black}
\SetBgContents{
        \begin{tikzpicture}[overlay,remember picture]
                \draw [line width=1,double, red!25] ($(current page.north west) + (2cm,-2cm) $) 
                rectangle
                ($(current page.south east) + (-2cm, 2cm) $);
                \node[anchor=south west] at ($(current page.north west) + (2cm,-2cm)$) {Kungl. Maj:ts Nåd. Proposition Nr 110.};
        \end{tikzpicture}
}


\begin{document}
{\centering

        {\calfamily
                {\Huge Sveriges Riksdag}\\
                {\Large Vid dess sammantrade,}\\ % I couldn't find a blackletter font with 'ä'
        }
        å Helgeandsholmen, Stockholm, mars månads 29e dag, Herrens år ett tusen
        nio hundrade och tolv.\\[0.5cm]

        \rule{5cm}{0.2ex}\\[0.5cm]
        {\large\scshape Proposition}\\
        För införandet i riksdagsordningen av kvinnlig rösträtt.\\[0.5cm]
        \rule{5cm}{0.2ex}\\[0.5cm]
}

{\itshape
        Under åberopande af bifogade i statsrådet förda protokoll vill Kungl.
        Maj:t härmed till Riksdagens Pröfning i grundlagsenlig ordning
        framlägga följande
}

Att lydelsen i riksdagsordningens §§ 9, 16, 19 och 21 ändras enl.

\section{9}
Till ledamöter i första kammaren kunna endast väljas män och kvinnor, som
uppnått  \rule[0.5ex]{2cm}{0.2ex} sin befattning.

\section{16}
Valrätt tillkommer enhvar välfrejdad svensk medborgare, såväl man som kvinna,
från och med kalenderåret näst efter det, hvarunder han uppnått tjugufyra års
ålder, dock ej
\begin{enumerate}[label=\alph*)]
        \item den som står under förmynderskap eller är i konkurstillstånd;
        \item gift kvinna, som ej vunnit boskillnad och hvilkens man är i
                konkurstillstånd;
        \item den som häftar för understöd, hvilket under löpande eller
                sistförflutna kalenderåret af fattigvårdssamhälle tilldelats
                den häftande själf eller dennes hustru eller minderåriga barn;
        \item gift kvinna, hvilkens man häftar för understöd, som under löpande
                eller sistförflutna kalenderåret af fattigvårdssamhälle
                tilldelats henne eller makarnas minderåriga barn;
        \item den som icke erlagt de honom påförda utskylder till stat och
                kommun, hvilka förfallit till betalning under de tre
                sistförflutna kalenderåren ;
        \item gift kvinna, som ej vunnit boskillnad och hvilkens man icke
                erlagt de honom påförda utskjdder till stat och kommun, hvilka
                förfallit till betalning under de tre sistförflutna
                kalenderåren;
        \item värnpliktig, som icke fullgjort de honom till och med utgången af
                sistförflutna kalenderåret åliggande värnpliktsöfningar.  
\end{enumerate}

Till efterrättelse vid val skall finnas röstlängd; och skall, på sätt i
vallagen finnes närmare bestämdt, valrätten grundas på förhållandena vid tiden
för röstlängdens tillkomst, ändå att förändring före valet inträffar.

\section{19}
Till ledamöter i andra kammaren kunna endast utses män och kvinnor, som äga
valrätt inom valkretsen eller, där fråga är om stad, bestående af flera
valkretsar, inom någon af dessa.

\section{21}
Riksdagsman \rule[0.5ex]{2cm}{0.2ex} tre lagtima riksdagar.  Kvinna är
berättigad att afsäga sig riksdagsmannauppdrag, äfven om icke något af nu
nämnda skäl är för handen.

Afsägelse \rule[0.5ex]{2cm}{0.2ex} Konungens befallningshafvande.

Kungl. Maj:t förblifver Riksdagen med all Kungl. nåd och ynnest städse
välbevågen.

\vspace{2cm}
{\centering\Huge\scshape
        Gustaf.\\
}

{\raggedleft\itshape
        Gust. Sandström.\\
}

\cofeAm{0.5}{1}{50}{0cm}{-13cm}

\end{document}
